\section{Related Work}
\label{sec:related}

\paragraph{Benchmarks for LLM agents in security operations} ExCyTIn-Bench~\citep{excytin}, SIABench~\citep{siabench}, and the Cyber Defense Benchmark~\citep{cyberdefensebench} evaluate agents on threat investigation, incident analysis, and threat hunting; AuditBench evaluates them on audit-log investigation tasks such as alert triage and persistence hunting~\citep{auditbench}, DiagChain and CyberClear on reconstructing attack chains from telemetry~\citep{diagchain,cyberclear}, and Cybench on offensive tasks~\citep{cybench}. SIR-Bench~\citep{sirbench} is closest in motivation: it scores whether an incident-response agent discovers new evidence rather than repeating the alert. These works build benchmarks; we audit existing ones with investigation-free baselines and per-question logs, and our checks apply to new benchmarks as well. Systems such as CORTEX, Clouseau, and SENTINEL-RL build triage and investigation agents~\citep{cortex,clouseau,sentinelrl}; claims about such systems rest on benchmarks of the kind we audit.

\paragraph{Re-evaluations in security} Arp et al.\ catalog pitfalls of machine learning in security, including spurious correlations and lab-only evaluation~\citep{arp2022dos}. Engelen et al.\ find labeling and generation errors in CIC-IDS2017~\citep{engelen2021cicids}, the source of SIABench's triage alerts. Bilot et al.\ re-evaluate provenance-based intrusion detectors and show that simple baselines match complex systems~\citep{bilot2025simpler}. We bring this kind of re-evaluation to LLM agent benchmarks, where the shortcut can sit in the question text rather than in features, and where per-question agent logs let us test whether the models use it.

\paragraph{Shortcuts and artifacts in machine learning} Shortcut learning~\citep{geirhos2020shortcut} and annotation artifacts~\citep{gururangan2018artifacts} describe models that exploit cues in a benchmark. Our ExCyTIn result is the converse case: the cue exists, the models do not use it, and the benchmark still fails to measure what it intends, because success depends only weakly on the investigation it requires.
